\paragraph{Why the evaluation object transfers beyond EEG.}
The general issue is attribution under unequal resources. If a component-bearing
system is scored as $S(1,R_1;T)$ and its comparator as $S(0,R_0;T)$, their
difference combines the component, the change from $R_0$ to $R_1$, and their
interaction. A component-specific increment instead compares $S(1,R;T)$ with
$S(0,R;T)$ on fixed support $T$. Related evidence shows why this distinction is
material: higher multimodal performance need not identify cross-modal
interaction \citep{hessel2020interactions}; controlled vision--language
comparisons attribute many reported differences to training data and
hyperparameters \citep{bugliarello2021unmasked}; and active learning, domain
adaptation, and personalized federated learning change with annotation or local
data budgets \citep{hacohen2022budget,bai2021pretrain,marfoq2022local}.
The invariant depends on the claim: an active-learning study may treat selected
labels as the intervention and hold annotation cost fixed, whereas a modality
ablation that claims added information should share calibration labels. Our
cross-domain contribution is therefore a reusable resource-contract test for
component attribution. FACED validates its consequence for EEG; it does not
estimate the prevalence of mismatch in other fields.
